\chapter{Requisitos y Diseño de la Solución}
\label{ch:requisitos-diseno}

\section{Requisitos funcionales y no funcionales}

\subsection{Requisitos funcionales}

Los requisitos siguientes describen el comportamiento esperado del sistema desde la perspectiva
del usuario y del operador.

\begin{enumerate}
  \item \textbf{Interfaz de acceso.} El sistema debe exponer todas sus funcionalidades a través
        de una \ac{API} REST como único punto de entrada programático. Las operaciones de ingesta,
        consulta y revisión deben ser accesibles sin intervención de la interfaz gráfica.

  \item \textbf{Ingesta de reuniones.} El sistema debe aceptar grabaciones de audio en formato
        MP3, WAV o M4A, así como transcripciones preformateadas en YAML o JSON. Ambas rutas deben
        converger en una representación interna común antes de continuar el procesamiento.

  \item \textbf{Extracción estructurada.} A partir de la transcripción, el sistema debe identificar
        nodos de cuatro tipos: decisiones (\texttt{decision}), riesgos (\texttt{risk}), preguntas
        abiertas (\texttt{open\_question}) e ítems de acción (\texttt{action\_item}). Cada nodo
        debe quedar anclado a una cita literal del texto fuente.

  \item \textbf{Resolución de relaciones.} El sistema debe inferir relaciones semánticas entre los
        nodos extraídos y los nodos preexistentes en la base de conocimiento, clasificándolas en
        siete tipos: \texttt{supersedes}, \texttt{amends}, \texttt{conflicts\_with},
        \texttt{depends\_on}, \texttt{mitigates}, \texttt{blocks} y \texttt{resolves}.

  \item \textbf{Puntuación de confianza.} Cada nodo extraído debe recibir una puntuación de
        confianza en $[0, 1]$ antes de la revisión humana. La puntuación combina señales
        heurísticas basadas en reglas y, opcionalmente, un agente de verificación independiente.

  \item \textbf{Revisión y corrección humana.} Los nodos con puntuación inferior al umbral
        configurado deben quedar en estado de revisión pendiente. Un revisor humano debe poder
        aprobar, rechazar o editar cada nodo antes de que éste entre en la base de conocimiento.
        Debe existir también un mecanismo de aprobación masiva por umbral.

  \item \textbf{Base de conocimiento consultable.} Los nodos aprobados deben persistir en una
        base de conocimiento con forma de grafo, consultable por tipo, estado, equipo, proyecto,
        dominio y rango de fechas. El sistema debe permitir búsqueda semántica y travesía del
        grafo para explorar el impacto de un nodo sobre otros.

  \item \textbf{Marco de evaluación.} El sistema debe incluir un marco de evaluación funcionalmente
        determinista con corpus etiquetado, que mida la calidad de extracción y bloquee el procesamiento
        de datos reales si no se alcanzan los objetivos de calidad definidos.

\end{enumerate}

\subsection{Requisitos no funcionales}

Los requisitos siguientes describen las propiedades de calidad que el sistema debe satisfacer con
independencia de la funcionalidad concreta que se esté ejecutando.

\begin{itemize}
  \item \textbf{Trazabilidad.} Cada nodo debe referenciar la cita literal del texto fuente que
        lo originó. Todos los prompts e inferencias deben quedar registrados para auditoría.

  \item \textbf{Seguridad.} El acceso a la \ac{API} debe estar controlado mediante claves de \ac{API} con
        modelo de roles. Las entradas de audio deben validarse por contenido (firma de bytes), no
        por extensión o cabecera. Las instrucciones del sistema deben estar aisladas del contenido
        procesado para prevenir la inyección de prompts.

  \item \textbf{Extensibilidad.} Los proveedores de LLM, transcripción, almacenamiento vectorial
        y secretos deben ser intercambiables sin modificar la lógica del \textit{pipeline}. Los
        cambios de proveedor deben ser una modificación de configuración, no una refactorización.

  \item \textbf{Observabilidad.} El sistema debe registrar el uso de tokens, la latencia por
        etapa, las métricas de confianza y los errores en MLflow~\cite{zaharia2018mlflow} de
        forma automática, sin instrumentación manual en los agentes.

  \item \textbf{Reproducibilidad.} El marco de evaluación debe producir resultados consistentes
        entre ejecuciones sobre el mismo corpus, dentro de los límites del no determinismo
        inherente a los \ac{LLM}.

  \item \textbf{Atomicidad.} Cada escritura en la base de conocimiento debe ser una transacción
        única. No puede existir estado parcial tras un fallo.

  \item \textbf{Resiliencia.} Las llamadas a proveedores externos —LLM, transcripción, reranking,
        almacenamiento— deben estar sujetas a reintentos con \textit{backoff} exponencial y
        \textit{jitter}, y a tiempos de espera máximos configurables. El consumo de tokens debe
        estar acotado mediante límites por llamada y por ejecución, tanto para LLM como para
        embeddings. El sistema debe continuar operando ante fallos transitorios sin requerir
        intervención manual.
\end{itemize}

% ---------------------------------------------------------------------------
\section{Arquitectura de la aplicación}

\textbf{Seshat} se organiza en cuatro capas con flujo de control estrictamente unidireccional,
tal como muestra la figura~\ref{fig:arquitectura}.

\begin{figure}[H]
  \centering
  \includegraphics[width=\textwidth]{architecture-diagram.png}
  \caption{Arquitectura de \textbf{Seshat}: la interfaz envía trabajos a la \ac{API} por HTTP; la
           \ac{API} los encola para el Worker; el Worker ejecuta el procesamiento y persiste los
           resultados en la capa de almacenamiento.}
  \label{fig:arquitectura}
\end{figure}

Cada capa tiene una responsabilidad delimitada, que se describe a continuación:

\begin{enumerate}
  \item \textbf{Interfaz de usuario}: envía trabajos, muestra su estado y presenta los nodos
        pendientes de revisión al operador. Se comunica exclusivamente con la \ac{API} por HTTP.
  \item \textbf{API}: autentica y valida las peticiones, registra el trabajo en la base de datos
        y lo encola para su procesamiento asíncrono.
  \item \textbf{Worker}: consume los trabajos de la cola y ejecuta el procesamiento
        correspondiente, persistiendo los resultados en la capa de almacenamiento.
  \item \textbf{Capa de almacenamiento}: Postgres para el estado operativo y la base de
        conocimiento; pgvector para los embeddings; almacenamiento de objetos compatible con S3
        para artefactos de audio, transcripciones y resultados de extracción.
\end{enumerate}

% ---------------------------------------------------------------------------
\section{Selección tecnológica}

\subsection{Lenguaje y frameworks}

\textbf{Seshat} está implementado en Python, el lenguaje dominante en el ecosistema de IA
generativa. Las principales dependencias son:

\begin{itemize}
  \item \textbf{FastAPI}: API REST asíncrona con validación de esquemas y documentación OpenAPI autogenerada.
  \item \textbf{LangChain}: orquestación de agentes LLM con soporte para múltiples proveedores; integra pgvector como almacén vectorial mediante \texttt{langchain-postgres}.
  \item \textbf{MLflow}: observabilidad: registro automático de trazas, uso de tokens, latencia y métricas de evaluación.
  \item \textbf{Pydantic}: validación de modelos de datos y configuración centralizada desde variables de entorno.
\end{itemize}

\subsection{Proveedores de modelos externos}

Todos los proveedores son intercambiables mediante configuración. El sistema distingue cuatro
roles:

\begin{itemize}
  \item \textbf{Transcripción:} AssemblyAI por defecto; sustituible por OpenAI Whisper u otros
        mediante la interfaz \texttt{AbstractTranscriber}.

  \item \textbf{LLM:} OpenAI, Anthropic, Azure OpenAI y Amazon Bedrock Converse, todos a través
        de la abstracción de LangChain.

  \item \textbf{Embeddings:} OpenAI, Azure OpenAI y Cohere sobre pgvector, otra vez mediante la
        abstracción de LangChain.

  \item \textbf{Reranking:} Cohere y Voyage como paso opcional de reordenación de candidatos
        RAG, encapsulados detrás de la interfaz \texttt{AbstractReranker}.
\end{itemize}

Un requisito de diseño relevante afecta al agente de verificación de confianza
(\textit{grounding}): debe usar un proveedor distinto al de los agentes de identificación, ya que
verificar con el mismo modelo produce errores correlacionados que anulan el valor de la
señal~\cite{zheng2023judging,cohen2023lmvslm}.

\subsection{Almacenamiento}

La capa de almacenamiento combina tres tecnologías con roles claramente separados:

\begin{itemize}
  \item \textbf{PostgreSQL}: base de datos principal; almacena el estado operativo (trabajos,
        claves de \ac{API}) y los nodos y relaciones de la base de conocimiento.
  \item \textbf{pgvector}: extensión de Postgres para búsqueda por similitud vectorial. Elimina
        la necesidad de un servicio vectorial separado y mantiene las escrituras en la base de
        conocimiento y los embeddings sincronizadas.
  \item \textbf{S3 / LocalStack}: almacenamiento de artefactos: audio original, transcripciones
        y resultados de extracción. LocalStack emula la \ac{API} de AWS en desarrollo.
\end{itemize}

La gestión de secretos (claves de \ac{API}, cadena de conexión de base de datos, etc) se delega en AWS
Secrets Manager, también emulado por LocalStack en desarrollo.

% ---------------------------------------------------------------------------
\section{Diseño de datos y pipelines}

\subsection{Modelo de datos: el grafo de conocimiento}

El conocimiento extraído de las reuniones se representa como un grafo de
conocimiento~\cite{hogan2021knowledge}: un grafo dirigido y etiquetado donde, los nodos pueden ser
entidades de cuatro tipos (ver tabla~\ref{tab:node-types}) y las aristas son relaciones semánticas
entre ellos de siete tipos, agrupados en dos categorías (ver tabla~\ref{tab:rel-types}). La dimensión
temporal es inherente al modelo: cuando una reunión posterior introduce un nodo que reemplaza o matiza
a uno anterior, la relación queda registrada explícitamente y el estado del nodo destino se actualiza
en consecuencia. El grafo preserva así la historia completa del conocimiento organizativo, no sólo su
estado actual.

\begin{figure}[H]
  \centering
  \includegraphics[width=\textwidth]{kg_example.drawio.png}
  \caption{Ejemplo de grafo de conocimiento en \textbf{Seshat}. El nodo con borde discontinuo
           representa una decisión previa en estado \texttt{SUPERSEDED}; las relaciones entre
           nodos de distinto tipo se muestran con flechas de línea discontinua.}
  \label{fig:kg-example}
\end{figure}

\begin{table}[H]
  \centering
  \begin{tabularx}{\textwidth}{l>{\raggedright\arraybackslash}X}
    \toprule
    \textbf{Tipo} & \textbf{Descripción} \\
    \midrule
    \texttt{decision}
      & Compromiso explícito y de ámbito grupal sobre qué opción adoptar, qué política seguir
        o qué restricción imponer. Tiene consecuencias operativas más allá de la conversación. \\
    \texttt{risk}
      & Modo de fallo concreto, bloqueante activo o incertidumbre con una consecuencia declarada
        que el grupo trata como sustantiva y no resuelta. \\
    \texttt{open\_question}
      & Pregunta cuya respuesta el grupo necesita pero aún no ha determinado; su resolución
        condiciona decisiones o trabajo futuro. \\
    \texttt{action\_item}
      & Tarea asignada a una persona o equipo con resultado esperado y, típicamente, fecha límite. \\
    \bottomrule
  \end{tabularx}
  \caption{Tipos de nodo en el grafo de conocimiento de \textbf{Seshat}.}
  \label{tab:node-types}
\end{table}

\begin{table}[H]
  \centering
  \begin{tabularx}{\textwidth}{ll>{\raggedright\arraybackslash}X}
    \toprule
    \textbf{Tipo} & \textbf{Categoría} & \textbf{Definición} \\
    \midrule
    \texttt{supersedes}      & Mismo tipo   & El nodo fuente reemplaza permanentemente al destino \\
    \texttt{amends}          & Mismo tipo   & El nodo fuente matiza o actualiza parcialmente al destino \\
    \texttt{conflicts\_with} & Mismo tipo   & Los dos nodos están en contradicción activa \\
    \texttt{depends\_on}     & Mismo tipo   & El nodo fuente presupone al destino \\
    \texttt{mitigates}       & Tipo cruzado & Un riesgo queda mitigado por una decisión \\
    \texttt{blocks}          & Tipo cruzado & Un riesgo bloquea una decisión o pregunta abierta \\
    \texttt{resolves}        & Tipo cruzado & Una decisión o tarea cierra una pregunta abierta \\
    \bottomrule
  \end{tabularx}
  \caption{Tipos de relación en el grafo de conocimiento de \textbf{Seshat}.}
  \label{tab:rel-types}
\end{table}

Cada nodo almacena su contenido —título, descripción y cita literal del texto fuente— junto con
una puntuación de confianza y metadatos de auditoría. El contenido es estable por diseño una vez
aprobado; el único campo que cambia automáticamente es el estado, actualizado cuando una reunión
posterior introduce una relación de supersesión o enmienda.

El grafo en su conjunto sigue un principio de \textit{append-and-state-only}: no existe operación
de borrado ni sobreescritura. El único matiz es que el contenido de un nodo aprobado puede
corregirse mediante una operación explícita reservada a operadores y administradores, que exige
justificación escrita y queda registrada en el historial de auditoría.

\subsection{Ciclo de vida de un trabajo}

Un trabajo \textit{pipeline} en \textbf{Seshat} representa el procesamiento completo de una reunión,
desde la recepción del audio o texto hasta la persistencia de los resultados en la base de
conocimiento. Recorre seis estados principales: \texttt{PENDING}, \texttt{TRANSCRIBING},
\texttt{IDENTIFYING}, \texttt{AWAITING\_REVIEW}, \texttt{RESOLVING} y \texttt{WRITING}, más
los estados terminales \texttt{DONE} y \texttt{FAILED}. Las dos fases de trabajo intensivo con
LLM —identificación de nodos y resolución de relaciones— están separadas por la revisión humana
opcional (\textit{Human-in-the-Loop}, HITL): si hay nodos pendientes de aprobación y el modo
automático no está activo, el \textit{pipeline} se detiene en \texttt{AWAITING\_REVIEW} hasta
que el operador interviene. Los fallos recuperables permiten reintentar el trabajo; los fatales
requieren un nuevo envío.

La figura~\ref{fig:job-lifecycle} muestra las transiciones posibles.

\begin{figure}[H]
  \centering
  \includegraphics[width=\textwidth]{job_lifecycle.drawio.png}
  \caption{Ciclo de vida de un \textit{pipeline} en \textbf{Seshat}. \texttt{IDENTIFYING} y
           \texttt{RESOLVING} corresponden a las dos pasadas de extracción; el paso
           \textit{HITL} indica la revisión humana opcional; y en \texttt{WRITING} se
           persisten los nodos y sus relaciones en la base de conocimiento.
           \linebreak
           Las flechas discontinuas rojas señalan transiciones a error y la flecha inferior
           muestra el reintento mediante llamada a la \ac{API}.}
  \label{fig:job-lifecycle}
\end{figure}
